\begin{theorem}[Markov’s Inequality]
Let \((\Omega,\mathcal{F},\mathbb{P})\) be a probability space and let
\(X:\Omega\to[0,\infty]\) be a non‑negative random variable.
For any real number \(a>0\),

\[
\mathbb{P}(X\ge a)\le\frac{\mathbb{E}[X]}{a}.
\]
\end{theorem}

\begin{proof}
\textbf{1. Indicator of the event \(\{X\ge a\}\).}  
Define the indicator (characteristic) function

\[
I_{\{X\ge a\}}(\omega)=
\begin{cases}
1, &\text{if } X(\omega)\ge a,\\[2mm]
0, &\text{if } X(\omega)< a .
\end{cases}
\]

By definition, \(I_{\{X\ge a\}}\) is a non‑negative random variable taking only the
values \(0\) and \(1\).

\medskip
\textbf{2. Pointwise inequality.}  
Because \(a>0\) and \(X(\omega)\ge0\) for every \(\omega\),

\begin{align*}
I_{\{X\ge a\}}(\omega)\le\frac{X(\omega)}{a}\qquad\text{for all }\omega\in\Omega.
\tag{1}
\end{align*}

\noindent\textit{Justification.}
\begin{itemize}
\item If \(X(\omega)\ge a\), then \(I_{\{X\ge a\}}(\omega)=1\) and
\(\frac{X(\omega)}{a}\ge\frac{a}{a}=1\), so \(1\le\frac{X(\omega)}{a}\).
\item If \(X(\omega)< a\), then \(I_{\{X\ge a\}}(\omega)=0\) while
\(\frac{X(\omega)}{a}\ge0\); thus \(0\le\frac{X(\omega)}{a}\).
\end{itemize}
Hence (1) holds in every case.

\medskip
\textbf{3. Monotonicity of expectation.}  
The expectation operator is monotone for non‑negative random variables:
if \(Y(\omega)\le Z(\omega)\) for all \(\omega\), then \(\mathbb{E}[Y]\le\mathbb{E}[Z]\).
Applying this to (1) gives

\begin{align*}
\mathbb{E}\!\big[\,I_{\{X\ge a\}}\,\big]\le
\mathbb{E}\!\bigg[\frac{X}{a}\bigg].
\tag{2}
\end{align*}

\medskip
\textbf{4. Linearity of expectation.}  
Since \(\frac{X}{a}= \frac{1}{a}\,X\) and \(\frac{1}{a}\) is a deterministic constant,

\begin{align*}
\mathbb{E}\!\bigg[\frac{X}{a}\bigg]=\frac{1}{a}\,\mathbb{E}[X].
\tag{3}
\end{align*}

\medskip
\textbf{5. Expectation of an indicator equals a probability.}  
For any event \(A\in\mathcal{F}\), \(\mathbb{E}[I_A]=\mathbb{P}(A)\); therefore

\begin{align*}
\mathbb{E}\!\big[\,I_{\{X\ge a\}}\,\big]=\mathbb{P}(X\ge a).
\tag{4}
\end{align*}

\medskip
\textbf{6. Combine the previous relations.}  
Substituting (3) and (4) into (2) yields

\begin{align*}
\mathbb{P}(X\ge a)\le\frac{\mathbb{E}[X]}{a},
\end{align*}
which is precisely Markov’s inequality.

\medskip
\textbf{7. Discussion of assumptions.}  
The argument uses only that \(X(\omega)\ge0\) for all \(\omega\) (so that
\(X/a\) is non‑negative) and that \(a>0\) (to avoid division by zero). No
further integrability conditions are required beyond the existence (possibly
infinite) of \(\mathbb{E}[X]\).

\end{proof}